
\section*{5. Relations entre @@M0@@ et @@M1@@}

\textbf{5.1.} L'exemple donné dans 4.8 n'exclut pas qu'on puisse peut-être trouver, pour un schéma @@M2@@ quelconque, un homomorphisme d'anneaux
@@M3@@
tel que pour tout Module localement libre @@M4@@ sur @@M5@@, @@M6@@ transforme la classe de Chern @@M7@@ de Exp. V en la classe de Chern cohomologique de SGA 5 VII. Bien entendu, ici @@M8@@ désigne un nombre premier distinct des caractéristiques résiduelles de @@M9@@. Comme l'anneau @@M10@@ est engendré par les @@M11@@, on voit en tous cas que @@M12@@ est uniquement déterminé par la condition d'être compatible avec les classes de Chern.

\textbf{5.2.} Utilisant l'homomorphisme de Chern
@@M13@@
on en conclut, par passage aux gradués associés respectivement à la filtration du premier membre et la filtration par le degré du second, un homomorphisme de groupes gradués (mais pas d'anneaux !)
@@M14@@
qui satisfait à la relation
@@M15@@
où @@M16@@ désigne la classe de Chern @@M17@@-adique. On en conclut, par division à partir des @@M18@@, un homomorphisme d'anneaux :
@@M19@@
transformant classe de Chern en classes de Chern; et tel que l'image du premier membre soit contenue dans celle de @@M20@@; par suite, la question de l'existence de @@M21@@ se résout affirmativement pourvu qu'on remplace @@M22@@ par son quotient par son sous-groupe de torsion. Le premier cas douteux pour définir les composantes @@M23@@ de l'hypothétique homomorphisme @@M24@@ est celui de @@M25@@, on n'arrive à priori qu'à définir @@M26@@. La réponse à la question d'existence de @@M27@@ serait négative si par exemple on pouvait trouver trois faisceaux inversibles @@M28@@ sur @@M29@@ satisfaisant la relation
@@M30@@
